% Line Atlas chapter draft.
% Suggested dependencies: amsmath, amssymb, siunitx, booktabs, tabularx,
% tikz, pgfplots, xcolor, hyperref.

\chapter{Practical Spectroscopy and Measurement}
\label{chap:practical-spectroscopy}
\chaptermark{Practical Spectroscopy and Measurement}

The most elegant plasma diagnostic is only as reliable as the chain that connects a
photon to a number in a table.  Raw exposures contain detector signatures, sky and
thermal backgrounds, atmospheric absorption, cosmic rays, geometric distortion, and
an instrumental line-spread function.  Reduction removes or models these effects but
also introduces correlations and choices.  The final spectrum is not simply the sky
spectrum with noise added; it is the output of an estimator.

This chapter presents a practical workflow that applies across long-slit, fiber,
echelle, and integral-field spectroscopy.  Instrument-specific pipelines remain
essential, but the underlying bookkeeping is common.  Useful extended discussions
include \citet{Gray2005}, \citet{Howell2006}, and the optimal extraction method of
\citet{Horne1986}.

\section{Design the observation around the inference}

An observing proposal should begin with the physical comparison that will be made,
not with a list of attractive transitions.  For each required line, record:
\begin{itemize}
 \item its observed wavelength range after redshift and expected velocity offsets;
 \item the continuum or background against which it will be measured;
 \item expected flux, surface brightness, width, and spatial extent;
 \item nearby blends, sky lines, telluric bands, detector gaps, and order edges;
 \item the angular and spectral resolution required to distinguish the alternatives;
 \item whether the key observable is integrated flux, equivalent width, profile
 shape, centroid, doublet ratio, or spatial morphology.
\end{itemize}

The resolving power is
\begin{equation}
 R\equiv\frac{\lambda}{\Delta\lambda},
 \qquad \Delta v\simeq\frac{c}{R},
 \label{eq:resolving-power}
\end{equation}
where $\Delta\lambda$ and $\Delta v$ must refer to a stated width convention, often
the FWHM of the line-spread function (LSF).  Pixel sampling and resolution are not the
same.  Two or more pixels per FWHM avoid gross undersampling, but additional sampling
does not recover information absent from the optical LSF.

For an unresolved line measured in background-limited data, signal-to-noise improves
approximately as the square root of exposure time.  Real observations can depart from
this scaling because of read noise, residual flat-field structure, time-variable sky,
correlated resampling noise, and systematic background subtraction.  Exposure-time
calculators should be tested at the surface brightness and extraction aperture of the
science target rather than only for a point-source continuum.

\subsection{Worked example: can a spectrograph separate two components?}

Suppose two narrow components differ by $80\,\mathrm{km\,s^{-1}}$.  A spectrograph
with $R=3000$ has an instrumental FWHM near $c/R\simeq100\,\mathrm{km\,s^{-1}}$.
The components may broaden or skew the feature, but will not be cleanly resolved.
At $R=10{,}000$, the nominal FWHM is $30\,\mathrm{km\,s^{-1}}$, giving meaningful
separation if the signal-to-noise is adequate.  This estimate is not a guarantee:
the actual LSF can have wavelength-dependent wings, and two components with unequal
flux are harder to identify than equal components.

\section{Calibration frames and detector-level modeling}

A simplified raw detector model is
\begin{equation}
 D_{ij}=g^{-1}\left[S_{ij}+B_{ij}+K_{ij}+d_{ij}t\right]
 +b_{ij}+\epsilon_{ij},
 \label{eq:detector-model}
\end{equation}
where $S$ is source photocharge, $B$ sky or thermal background, $K$ scattered light,
$g$ gain, $b$ electronic bias, $d$ dark current, and $\epsilon$ stochastic read and
photon noise.  Nonlinearity, persistence, interpixel capacitance, charge transfer,
and bad pixels may require additional terms.

Bias or reference-pixel correction removes the electronic zero level.  Dark frames
estimate thermally generated signal and sometimes hot-pixel behavior.  Flat fields
measure pixel response and, depending on the instrument, illumination or blaze
structure.  Arc lamps, etalons, sky lines, or laser combs constrain the mapping from
detector position to wavelength.  Calibration frames are measurements with their own
noise and temporal mismatch; they should not be treated as exact arrays.

For CCD-like data in electrons, a basic variance model is
\begin{equation}
 V_{ij}=S_{ij}+B_{ij}+K_{ij}+d_{ij}t
       +n_{\rm read}\sigma_{\rm read}^2,
 \label{eq:variance-model}
\end{equation}
with appropriate gain factors if working in analog-to-digital units.  Operations such
as flat-field division, interpolation, and sky modeling must propagate variance.
Cosmic-ray rejection should mark altered samples; replacing a cosmic ray with a local
interpolation and assigning ordinary variance falsely creates information.

\section{Tracing, wavelength calibration, and rectification}

Spectral traces curve across detectors and may vary with exposure, fiber, slit
position, or echelle order.  A trace model should be smooth enough to follow optics
without following pixel noise.  Wavelength calibration fits known features with a
two-dimensional function $\lambda(x,y)$ or an instrument model.  Precision should be
evaluated using held-out or independent lines, not merely the residuals of lines used
in the fit.

Air and vacuum wavelengths must never be mixed silently.  Air's refractive index
shifts optical wavelengths enough to matter for velocity work.  The reference frame
must also be stated: topocentric, barycentric, heliocentric, local standard of rest,
or another convention.  A sign error in the observatory velocity correction can
exceed the physical signal being sought.

Several velocity definitions coexist.  Relative to a rest frequency $\nu_0$ or
wavelength $\lambda_0$,
\begin{align}
 v_{\rm radio} &= c\left(1-\frac{\nu}{\nu_0}\right),\\
 v_{\rm optical} &= c\left(\frac{\lambda-\lambda_0}{\lambda_0}\right),\\
 \beta_{\rm rel} &=
 \frac{(\lambda/\lambda_0)^2-1}{(\lambda/\lambda_0)^2+1}.
 \label{eq:velocity-conventions}
\end{align}
The first two are approximations that differ at non-negligible redshift.  A catalog
column named ``velocity'' is incomplete without a frame and convention.

Rectification places spectra on a convenient grid but correlates neighboring pixels.
Repeated interpolation broadens the effective LSF and compounds covariance.  A good
pipeline maps samples to the final grid once where practical, or retains native pixel
coordinates for forward modeling.

\section{Sky, thermal background, and telluric transmission}

Night-sky emission is structured in wavelength and time.  Near-infrared OH lines can
be much brighter than the target; small wavelength or LSF mismatches then leave
positive-negative residuals.  Sky subtraction is a model of the sky observed through
the same instrument, not a scalar offset.  Nodding, dedicated sky fibers, off-source
slit regions, or principal-component models each encode assumptions about spatial and
temporal variation.

Atmospheric transmission multiplies the extraterrestrial spectrum.  A telluric
standard or synthetic atmospheric model can estimate this transmission, but division
near saturated bands amplifies noise and cannot reconstruct photons that were not
transmitted.  Regions of very low transmission should usually be masked or fitted
jointly with a telluric model.  The atmospheric model and science model should be
convolved to the same LSF before comparison.

Thermal emission dominates at sufficiently long wavelengths and varies with telescope
and atmosphere temperature.  Background subtraction can be the largest term in the
error budget for faint extended emission.  A local background region may remove real
source emission if the object fills the field, while a remote sky exposure may not
match the instantaneous background.

\section{One-dimensional extraction}

For detector row $i$ and spatial sample $j$, let $D_{ij}$ be background-subtracted
counts, $P_{ij}$ a normalized spatial profile, and $V_{ij}$ variance.  The
inverse-variance optimal estimator is
\begin{equation}
 \widehat{f}_i=\frac{\sum_j P_{ij}D_{ij}/V_{ij}}
                       {\sum_jP_{ij}^2/V_{ij}},
 \qquad
 \mathrm{Var}(\widehat{f}_i)=
 \left(\sum_j\frac{P_{ij}^2}{V_{ij}}\right)^{-1}
 \label{eq:optimal-extraction}
\end{equation}
under the model assumptions \citep{Horne1986}.  Compared with an unweighted aperture,
optimal extraction downweights noisy spatial pixels.  It can be biased if the profile
is wrong, the object is extended, differential atmospheric refraction moves the
centroid, or emission-line morphology differs from continuum morphology.

Extraction aperture is part of the observable.  A narrow aperture can maximize
point-source signal-to-noise but exclude diffuse line halos.  Comparing line fluxes
from different seeing, slit widths, or instruments requires an aperture or throughput
model.  For spatially resolved targets it is often preferable to preserve the data
cube and define physical apertures during analysis.

\section{Flux calibration and response}

A spectrophotometric standard constrains the sensitivity function that converts count
rate to physical flux density.  The response contains telescope area, slit or fiber
throughput, optics, detector efficiency, atmospheric extinction, and sometimes an
aperture correction.  A smooth response model should not fit away real standard-star
absorption lines or imprint them on the target.

Absolute calibration and relative line-ratio calibration are different requirements.
Lines close in wavelength can have an accurate ratio even when absolute throughput is
uncertain.  Widely separated lines require accurate response shape and often an
attenuation correction.  Simultaneous photometry can rescale spectra affected by
clouds or slit loss, but a single broadband point cannot repair wavelength-dependent
losses.

For a source with flux density $f_\lambda$ and response $R(\lambda)$, detected
electrons scale as
\begin{equation}
 N_e\propto tA\int \frac{\lambda}{hc}
 f_\lambda(\lambda)R(\lambda)\,\mathrm{d}\lambda.
 \label{eq:photon-counting-response}
\end{equation}
The factor $\lambda/(hc)$ converts radiant energy to photon number.  Confusing
energy-weighted and photon-weighted response definitions produces color terms.

\section{A transparent analysis graph}

\begin{figure}[tbp]
\centering
\begin{tikzpicture}[
 node distance=7mm and 12mm,
 every node/.style={font=\small},
 box/.style={draw,rounded corners,align=center,minimum width=28mm,
 minimum height=9mm,fill=blue!7},
 warn/.style={draw,rounded corners,align=center,minimum width=28mm,
 minimum height=9mm,fill=orange!14},
 arrow/.style={->,thick}
]
 \node[box] (raw) {raw exposures\\and calibrations};
 \node[box,right=of raw] (det) {detector correction\\with masks/variance};
 \node[box,right=of det] (map) {trace and wavelength\\solution};
 \node[box,below=of map] (bg) {sky/telluric\\forward model};
 \node[box,left=of bg] (ext) {extraction or\\cube reconstruction};
 \node[box,left=of ext] (flux) {flux calibration\\and LSF};
 \node[box,below=of flux] (fit) {continuum plus\\line inference};
 \node[box,right=of fit] (phys) {physical model\\and diagnostics};
 \node[warn,right=of phys] (audit) {validation, provenance,\\posterior predictive checks};
 \draw[arrow] (raw)--(det); \draw[arrow] (det)--(map);
 \draw[arrow] (map)--(bg); \draw[arrow] (bg)--(ext);
 \draw[arrow] (ext)--(flux); \draw[arrow] (flux)--(fit);
 \draw[arrow] (fit)--(phys); \draw[arrow] (phys)--(audit);
 \draw[arrow,dashed] (audit.south) to[bend left=24]
   node[midway,below] {revise assumptions} (fit.south);
\end{tikzpicture}
\caption{Original reduction-to-inference graph.  Masks, variance or covariance, LSF,
and provenance flow with the spectrum.  Validation is a feedback loop rather than a
final cosmetic step.}
\label{fig:analysis-graph}
\end{figure}

Figure~\ref{fig:analysis-graph} emphasizes that every product should retain auxiliary
information: data-quality masks, propagated uncertainty, wavelength and spatial
coordinates, spectral resolution, exposure metadata, and links to calibration
versions.  A flux array without these is not analysis-ready.

\section{Continuum modeling}

Line measurements are conditional on a continuum model.  Possible components include
stellar populations and absorption, a nonthermal power law, nebular continuum, dust
emission, Fe complexes, scattered light, and residual response structure.  Local
polynomials are effective over narrow intervals but can absorb broad wings.  Physical
templates extrapolate more meaningfully but carry population and abundance mismatch.

The continuum and line should be fitted jointly whenever their uncertainties are
coupled.  Subtracting a best-fit continuum and then treating it as exact understates
line-flux uncertainty.  For Balmer emission, underlying stellar absorption can be
comparable to the nebular equivalent width, especially in weakly star-forming regions.
A flexible stellar-continuum model is then not optional.

Continuum placement can be stress-tested by changing windows, polynomial order, or
template library and examining the shift in line parameters.  These shifts are a
model-systematic term, not a reason to select whichever configuration yields the
smallest formal error.

\section{Forward line-profile fitting}

A common model for a local spectral interval is
\begin{equation}
 m(\lambda)=C(\lambda;\boldsymbol{c})+
 \left[\sum_{k=1}^{K}F_k\,
 p_k(\lambda;v_k,\sigma_k)\right]*L(\lambda),
 \label{eq:spectral-forward-model}
\end{equation}
sampled and integrated over the same pixels as the data.  Here each $p_k$ is normalized
to unit area, so $F_k$ is integrated flux.  Multiple transitions may share velocity
and dispersion; branching ratios may be fixed only when atomic physics, opacity, and
blending justify doing so.

With residual vector $\boldsymbol{r}=\boldsymbol{d}-\boldsymbol{m}$ and covariance
$\boldsymbol{C}$, the Gaussian log likelihood is
\begin{equation}
 -2\ln\mathcal{L}=
 \boldsymbol{r}^{\mathsf T}\boldsymbol{C}^{-1}\boldsymbol{r}
 +\ln|\boldsymbol{C}|+\mathrm{constant}.
 \label{eq:gaussian-correlated-likelihood}
\end{equation}
Ignoring off-diagonal covariance after resampling can underestimate uncertainties and
make broad low-amplitude residuals appear significant.  When the full covariance is
too costly, an empirical correction from blank-sky apertures or block bootstrap is
better than pretending pixels are independent.

For a Gaussian line in wavelength,
\begin{equation}
 p(\lambda)=\frac{1}{\sqrt{2\pi}\sigma_\lambda}
 \exp\!\left[-\frac{(\lambda-\lambda_c)^2}{2\sigma_\lambda^2}\right],
 \quad
 \mathrm{FWHM}=2\sqrt{2\ln2}\,\sigma_\lambda.
 \label{eq:gaussian-line-model}
\end{equation}
If both source and LSF are Gaussian,
$\sigma_{\rm obs}^2=\sigma_{\rm int}^2+\sigma_{\rm inst}^2$.  Subtraction in
quadrature is invalid for non-Gaussian LSFs, unresolved lines near the boundary
$\sigma_{\rm int}=0$, or blends.  Forward convolution handles these cases more
honestly.

\subsection{Model complexity and components}

Adding Gaussian components will almost always improve least-squares residuals.  A
component should be retained because it is reproducible, physically interpretable,
or required by predictive checks, not merely because an optimizer found it.  Compare
models using held-out pixels, information criteria with their assumptions stated, or
Bayesian evidence when priors are defensible.  Inspect whether inferred components
persist across independent transitions and spatial bins.

Parameters become non-identifiable when components are closer than the resolution or
one has low contrast.  In that regime, the data may constrain total flux and centroid
but not the individual widths and velocities.  Publishing a covariance or posterior
sample makes this distinction visible.

\section{Non-parametric profile measurements}

Non-parametric velocities can summarize asymmetric profiles without assigning each
feature to a Gaussian.  From the continuum-subtracted cumulative flux
\begin{equation}
 \mathcal{F}(v)=\frac{\int_{v_{\min}}^v F(v')\,\mathrm{d}v'}
 {\int_{v_{\min}}^{v_{\max}}F(v')\,\mathrm{d}v'},
 \label{eq:cumulative-line-flux}
\end{equation}
define $v_p$ by $\mathcal{F}(v_p)=p/100$.  Measures such as
$W_{80}=v_{90}-v_{10}$ and $v_{50}$ are common.  They depend on the integration
window, continuum, noise, and LSF.  Negative pixels and absorption complicate the
cumulative distribution, so the exact algorithm must be documented.

The flux-weighted centroid and second moment are
\begin{equation}
 \bar v=\frac{\int vF(v)\,\mathrm{d}v}{\int F(v)\,\mathrm{d}v},
 \qquad
 \sigma_v^2=\frac{\int(v-\bar v)^2F(v)\,\mathrm{d}v}{\int F(v)\,\mathrm{d}v}.
 \label{eq:line-moments}
\end{equation}
Moments are sensitive to weak wings and window choices.  Monte Carlo propagation
should perturb the continuum and correlated noise, not only individual flux pixels.

\section{Equivalent width, flux, and luminosity}

Rest-frame equivalent width separates a line from the local continuum scale.  For
emission, some communities use positive and others negative signs; the convention
must be explicit.  The line flux from a fitted profile is preferable to summing a
fixed window when blends or masked pixels matter.  Direct integration remains a
valuable model check.

Uncertainty in integrated flux for weights $w_i$ is
\begin{equation}
 \mathrm{Var}(F)=\boldsymbol{w}^{\mathsf T}
 \boldsymbol{C}\boldsymbol{w},
 \label{eq:integrated-flux-covariance}
\end{equation}
not merely $\sum_i w_i^2\sigma_i^2$ unless the covariance is diagonal.  Luminosity
then requires a cosmology and luminosity distance.  For gravitationally lensed
sources, quote whether luminosity is apparent or corrected by a magnification model.

\section{Detection significance and upper limits}

A line ``detected at $5\sigma$'' is incomplete unless the search procedure is stated.
If wavelength and width were scanned over many possibilities, the false-alarm rate
is larger than the single-trial Gaussian tail.  A matched-filter amplitude for a
known template $\boldsymbol{p}$ is
\begin{equation}
 \widehat{A}=\frac{\boldsymbol{p}^{\mathsf T}
 \boldsymbol{C}^{-1}\boldsymbol{d}}
 {\boldsymbol{p}^{\mathsf T}\boldsymbol{C}^{-1}\boldsymbol{p}},
 \qquad
 \sigma_A^2=\left(\boldsymbol{p}^{\mathsf T}
 \boldsymbol{C}^{-1}\boldsymbol{p}\right)^{-1}.
 \label{eq:matched-filter}
\end{equation}
The template includes the assumed intrinsic width and LSF.  Searching a template bank
requires calibrating the maximum statistic using blank spectra, permutations, or
simulations that reproduce the entire reduction.

An upper limit is conditional on line width, centroid window, profile, and confidence
construction.  For a nondetection of an unresolved line, a statement such as
``$3\sigma$ flux limit for a Gaussian with FWHM
$100\,\mathrm{km\,s^{-1}}$ at the systemic redshift'' is reusable.  A bare number is
not.  Negative fitted fluxes should generally remain in ensemble analyses; clipping
them to zero biases stacked and population estimates.

\subsection{Worked example: width-dependent flux limit}

Assume independent spectral samples of width $\delta v=20\,\mathrm{km\,s^{-1}}$
with flux-density uncertainty $s$.  A top-hat line of width
$\Delta v=100\,\mathrm{km\,s^{-1}}$ spans $N=5$ samples, so the integrated-noise
scale is $s\,\delta v\sqrt{N}$.  If instead one assumes
$\Delta v=400\,\mathrm{km\,s^{-1}}$, $N=20$ and the flux limit doubles.  Correlated
noise changes the square-root scaling.  This is why line-width assumptions belong in
upper-limit tables.

\section{Integral-field spectroscopy}

An integral-field unit samples two spatial dimensions and one spectral dimension, but
the reconstructed cube has a three-dimensional response: a spatial point-spread
function (PSF), spectral LSF, pixel footprint, and covariance.  Resampling detector
pixels into spaxels correlates both neighboring wavelengths and neighboring spatial
samples.

Common spatial binning methods target a continuum signal-to-noise, yet emission-line
signal-to-noise may follow a different morphology.  Binning based on the line being
tested can also introduce selection bias.  A defensible strategy defines bins using
an independent quantity, fits all spaxels hierarchically, or demonstrates through
injection tests that adaptive selection does not fabricate gradients.

Beam smearing mixes velocities within the PSF and can mimic high dispersion where a
velocity gradient is steep.  The observed cube should be compared to a source model
convolved spatially and spectrally.  Subtracting a scalar instrumental width from each
spaxel cannot correct spatial beam smearing.

\section{Validation by injection and repetition}

End-to-end injection tests place synthetic sources into raw or minimally processed
data, run the normal pipeline, and compare recovered parameters with truth.  They can
measure completeness, flux bias, centroid bias near sky lines, component
recoverability, and the effect of covariance.  Injections should span morphology,
line width, equivalent width, and background conditions rather than a single ideal
case.

Other high-value checks include:
\begin{itemize}
 \item fitting independent exposures and detector subsets separately;
 \item comparing integrated flux with direct aperture sums;
 \item verifying fixed atomic branching ratios where opacity is negligible;
 \item checking skylines or calibration lamps for LSF and wavelength residuals;
 \item measuring blank apertures with the same extraction and fitting procedure;
 \item examining posterior predictive residuals in detector and physical coordinates;
 \item repeating measurements under plausible continuum and mask choices.
\end{itemize}

\section{A minimum reproducible line record}

Every published line measurement should be accompanied, directly or in machine-readable
metadata, by:
\begin{enumerate}
 \item transition identifier, rest wavelength, air/vacuum convention, and atomic-data
 source;
 \item observed frame and velocity convention;
 \item aperture, PSF, slit losses, and extraction method;
 \item continuum model, fit window, masks, and blend constraints;
 \item LSF model and spectral sampling;
 \item flux or equivalent width with statistical and systematic uncertainties;
 \item covariance or a documented approximation to it;
 \item attenuation, lensing, and aperture corrections;
 \item nondetection assumptions and search trials;
 \item software version, configuration, calibration version, and provenance links.
\end{enumerate}
This record turns a table of numbers into scientific evidence that another researcher
can audit, combine, or reinterpret.
